% ============================================================================
%  Section 7: Discussion
%  Sources: blueprint secs. 2, 5-7; Sections 4-6 of this paper; FROT 2025_sub/12_app_quant.tex
%  (Construction of Outcomes; Log Preferences Estimation); tables/msm_lnc.tex; tables/D_schedules.tex;
%  tables/est_parameters.tex; fix/verified_facts.md (literature facts).
%  Derived quantities (exact arithmetic on reported table entries) are marked with % Derived comments.
% ============================================================================
\section{Discussion}\label{sec:discussion}

This section examines which of our results depend on measurement conventions and modeling choices.

% ---------------------------------------------------------------------------
\subsection{Converting Welfare into Dollars}\label{sec:disc_dollars}

The WTP for a reform is the proportional change in consumption, at every age from 30 to $T$ and in every state, that the cohort values as much as the reform (equation~\eqref{eq:ce}). Equation~\eqref{eq:mvpf} converts it into dollars with the present value of baseline per capita consumption from age 62, $\bar C^{0}_{62}$, for the reforms that start at 62, and from age 30, $\bar C^{0}_{30}$, for the savings subsidies. Only the lifetime base values the consumption change that defines the CE: on it, a transfer that raises consumption in the same proportion at every age and in every state, and induces no behavioral response, has an MVPF of one, whereas on the post-62 base its MVPF is $\bar C^{0}_{62}/\bar C^{0}_{30}$. The tables pin down this ratio. Every MVPF computed with the net cost equals the ratio of the WTP and Net entries times a factor that depends on baseline aggregates but not on the reform; computed from the rounded entries, the factor lies between 0.38 and 0.40 for the reforms that start at 62 and between 1.87 and 1.99 for the savings subsidies (Section~\ref{sec:fw_mvpf}). Apart from the Medicare reforms, which use their own baseline, every Net entry is a percentage change of the same baseline aggregate, net government spending in the status quo, so the two factors differ only through the consumption base. The lifetime base, which adds the consumption of the 32 years from age 30 to 61, is therefore about five times the post-62 base: the ratio lies between 4.7 and 5.2. % Derived: 1.87/0.40 = 4.7 and 1.99/0.38 = 5.2 (factors from Section 4.4). AUTHORS: with the inferred bases the factor is C0_s/|G0|; not printed (decision D2).

Three kinds of statements do not depend on the base: whether an MVPF is infinite, which depends only on the signs of WTP and net cost; the ranking of reforms within a family, which share a base; and the ranking across the families that start at 62 (Section~\ref{sec:results_map}), which share the post-62 base. The central results of the paper are of these kinds. The base matters for the levels of finite MVPFs and for two comparisons. First, on the post-62 base, the lower bounds of the savings MVPFs, 0.84 to 0.96, would be about one-fifth as large, roughly 0.16 to 0.20: close in level to the MVPFs of the NRA increases and below those of the old-age income-tax cuts. Their upper bounds would be about 4 for K1 and K2 and about 1 for K3. % Derived: reported bounds x (0.38 to 0.40)/(1.87 to 1.99): 0.84/5.2 = 0.16, 0.96/4.7 = 0.20; 20.37/5 = 4.1, 19.38/5 = 3.9, 5.29/5 = 1.1.
Where the subsidies rank among the reforms that start at 62 thus depends on how much of the indirect saving from taxing interest income, an upper bound in the model, materializes.

Second, on the lifetime base, every finite MVPF of the reforms that start at 62 would be about five times larger. The MVPFs of the old-age income-tax cuts and of every expansion of the benefit formula would exceed one: about 3.2 to 5.9 for P1.2--P1.9, 1.4 to 3.2 for the level reforms, 1.4 to 2.9 for the P2 and P3 reforms, and 1.1 to 1.4 for the income-tax cuts. % Derived: reported MVPFs x (4.7 to 5.2): 0.69 x 4.7 = 3.24, 1.14 x 5.2 = 5.93; 0.30 x 4.7 = 1.41, 0.61 x 5.2 = 3.17; 0.29 x 4.7 = 1.36, 0.56 x 5.2 = 2.91; 0.23 x 4.7 = 1.08, 0.26 x 5.2 = 1.35.
The MVPFs of the Medicare expansions, about 0.7 to 0.8, would lie below the point estimate of about 1.6 for the introduction of Medicare but inside its 95 percent confidence interval of 0.48 to 3.55; on the post-62 base, they lie below that interval (Section~\ref{sec:results_medicare}). These levels embed the utilitarian weights of our ex ante WTP across individuals who differ at age 30 and the low old-age wealth of the model (Section~\ref{sec:estimation}); they are not comparable in a strict sense to MVPFs computed from beneficiaries' own willingness to pay. Each dollar saved by raising the NRA would cost about 75 to 94 cents of willingness to pay, rather than the 16 to 18 cents measured on the post-62 base. % Derived: reported MVPFs x (4.7 to 5.2): 0.146 x 4.7 = 0.69, 0.155 x 5.2 = 0.81; 0.16 x 4.7 = 0.75, 0.18 x 5.2 = 0.94; 0.23 x 4.7 = 1.08; 0.29 x 4.7 = 1.36 (smallest finite benefit-formula MVPF, tables/tab_generosity_app_body.tex). Medicare CI: fix/verified_facts.md item 7.
A comparison of an MVPF with one is therefore meaningful only on the lifetime base.

% ---------------------------------------------------------------------------
\subsection{General Equilibrium}\label{sec:disc_ge}

Every reform is evaluated at a fixed wage per efficiency unit and a fixed interest rate, without balancing the budget. Fixed prices are exact in a small open economy, in which, with constant returns to scale, the world return to capital also fixes the wage. In a closed economy, a reform applied to every cohort would also move the wage, through aggregate labor input, the return to capital, through aggregate saving, and the taxes or spending that absorb its effect on the budget. Overlapping-generations analyses of Social Security reform incorporate these channels \citep{Conesa1999social,nishiyama2007does,Imrohoroglu2012,Kitao2014sustainable}.

The price channels are pecuniary. By the envelope theorem, individuals' responses to a small price change have no first-order effect on their welfare, which changes only through the direct effect of prices on their income. With constant returns to scale, the fall in the wage and the rise in the return that follow an increase in labor input leave aggregate pretax income unchanged at the initial factor supplies. They redistribute it from workers to asset holders, and they move the WTP of a reform only through differences in marginal utility between those who gain and those who lose. Their effect on the budget differs: additional labor raises output and the tax base in general equilibrium as in partial equilibrium, and the shift of income from labor to capital changes revenue by the difference between the tax rates on the two.

The sign of the fiscal saving of the reward-delay reforms is likely to survive these feedbacks. Its first source, the fall in retirement outlays, comes from behavioral responses within Social Security; for the claiming responses among them, the fiscal sign depends on prices only through the discount rate of Lemma~\ref{lem:claiming}: both terms of equation~\eqref{eq:claimnpv} are proportional to the PIA, so the wage affects their size but not their sign. Its second source is additional work. These reforms raise participation at 65--69 and, in their moderate versions, in every age group we track; a lower wage would reduce the revenue from each additional hour but, for a small reform, would not reverse the response that generates it. Returned through any instrument with a positive MVPF, the saving raises welfare further while balancing the budget (Appendix~\ref{app:derivations}). The saving of the NRA increases is mechanical, so its sign does not depend on prices either. Correcting the size of these effects would require, beyond the cohort's responses that the model delivers, the elasticities of the wage and the return to capital with respect to aggregate labor and saving, each reform's effect on these aggregates summed over the cohorts alive in a steady state, the tax rates on labor and capital income, and the instrument that balances the budget.

% ---------------------------------------------------------------------------
\subsection{Statutory Ages as Reference Points}\label{sec:disc_reference}

In our model, the NRA matters only through the retirement benefits, DI eligibility, and earnings-test rules it determines. Evidence suggests that statutory ages matter beyond these incentives. \citet{seibold2021reference} documents bunching of retirements at statutory ages in Germany, including ages at which the financial incentives to retire change little or not at all, and interprets it as evidence that statutory ages act as reference points. In Switzerland, an increase in women's full retirement age delayed both claiming and labor market exit by more than a standard life-cycle model predicts, whereas changes in the penalties for early claiming with the full retirement age held fixed had much smaller effects \citep{Lalive2023how}. In the United States, the spike in benefit claiming at the full retirement age moved with it as the 1983 Amendments raised it \citep{behagel2012}; when it rose from 65 to 66, however, labor force exit remained concentrated at 65 \citep{Deshpande2024how}.

Reference dependence would change the fiscal effect of an NRA increase in ways our model does not capture, and it would make its welfare cost harder to measure. If claims and exits both follow the NRA, as in the Swiss evidence, an increase raises participation below the new NRA by more than our model predicts; for a given WTP, the larger fiscal saving would lower its MVPF further. If claims follow the NRA but exits do not, as in the US evidence, the rise in participation at 65--69 that our model predicts, 0.60 to 3.27 points, may overstate the response. Because the indirect effects of the NRA increases are small (Section~\ref{sec:results_claiming}), their saving would then change mainly through the shift of claims toward the new NRA, whose effect on outlays depends on the increments of the schedule relative to claimants' actuarially fair rates (Section~\ref{sec:fw_claiming}). On the welfare side, if the reference point is part of preferences, moving it changes utility directly; if it is a norm followed without optimizing, the envelope argument under which behavioral responses enter the MVPF only through the budget (Section~\ref{sec:fw_decomp}) fails. In both cases, a WTP computed from preferences over consumption and leisure misses part of the welfare effect of the reform, and the behavioral evidence does not determine its sign.

The reward-delay reforms leave the NRA, and any reference point at it, at 65. Reference dependence would bear on them through the response to the credit: if the NRA anchors claiming, fewer individuals would delay beyond it when the credit rises. The Swiss finding that early-claiming penalties move behavior much less than the full retirement age itself points in this direction \citep{Lalive2023how}. Because the mechanical effect of a DRC increase on retirement outlays is a cost, its infinite MVPF rests on these responses (Section~\ref{sec:results_claiming}), and strong anchoring would reduce its saving. Our model matches the concentration of exits in the 60s (\autoref{fig:fit_labor}) but attributes it to incentives, health, and the fixed costs of work; to the extent that anchoring at 65 accounts for part of it, the model may overstate the response to the credit.

% ---------------------------------------------------------------------------
\subsection{Cohort Rules and Current Law}\label{sec:disc_cohort}

Our status quo represents the rules faced by the cohorts born between 1920 and 1940 with an NRA of 65, a reduction for early claiming of $6\tfrac{2}{3}$ percent of the PIA per year, and a DRC of 4.5 percent, approximately the cohort average of credits that rose with the year of birth from 3 to 7 percent. AN2 and AD1 correspond to the end points of the two changes to claiming incentives in the 1983 Amendments, which raised the NRA toward 67 and the DRC to 8 percent, both phased in by birth cohort (Section~\ref{sec:results_claiming}; \citealp{Svahn1983social}). Combined, the two end points give the current-law schedule, which pays 70 percent of the PIA at 62 and 124 percent at 70. % Derived: AN2 schedule (tables/D_schedules.tex) with an 8 percent DRC after 67: 100 + 3 x 8 = 124.
Our grid does not contain this combination, whose effects need not equal the sum of those of AN2 and AD1, because the response to a higher credit depends on the age from which it accrues.

The grid still speaks to the direction of further reforms. Because the CE is a log ratio and the Net entries of a table share their baseline aggregate, the difference between two rows evaluates the move from one reform to the other. Relative to the status quo, every credit from 8 to 10 percent in our grid pays for itself. Relative to an 8 percent credit at an NRA of 65 (AD1), a credit of 9, 9.5, or 10 percent (AD3--AD5) raises WTP by 0.07 to 0.14 percent and lowers net government spending by 0.21 to 0.33 percent of its status-quo value, whereas a credit of 8.5 percent (AD2) changes WTP and net spending by amounts within the precision of the simulations. % Derived: row differences in tables/tab_claiming_body.tex (WTP 0.23-0.16 to 0.30-0.16; Net -0.66+0.45 to -0.78+0.45; AD2: +0.03, +0.01).
The reforms AD1--AD5 and AL2--AL9 are at least as steep, at every claiming age, as the current-law schedule, which rises by 5 percent of the PIA per year from 62 to 64, by $6\tfrac{2}{3}$ percent from 64 to 67, and by 8 percent after 67 (Appendix~\ref{app:reforms}). % Derived: year-to-year increments in tables/D_schedules.tex; AD1: 6.67, 6.66, 6.67, 8, 8, 8, 8, 8; AL1 (7 after 67) is less steep.
For the cohorts we study, and relative to their schedule, steeper schedules at an NRA of 65 raised welfare and saved money (Section~\ref{sec:results_claiming}); because these schedules differ from current law in level as well as in slope, the grid does not indicate whether steepening the current-law schedule would do so.

The grid does not say whether such reforms pay for themselves from current law, and four features of our setting favor the finding that rewards for delay pay for themselves. First, the baseline credit is 4.5 percent, at which delay past 65 was, on average, not actuarially advantageous (Section~\ref{sec:results_claiming}); under current law the credit is already 8 percent and accrues from 67, when fewer individuals still work, and who would respond to a further increase can be determined only by solving the model under current law. Second, the horizon ends at 85, which truncates the stream of benefits that a delay raises. Third, the government discounts at the net return $r$. Fourth, the individuals in our model receive no spousal or survivor benefits. By equation~\eqref{eq:faircredit}, survival beyond 85 and a lower discount rate each lower the actuarially fair credit, and survivor benefits would do the same by lengthening the stream of benefits that a primary earner's delay raises; each lowers the saving from a given delay. Spousal benefits work differently: they earn no delayed retirement credit, so a higher credit would reach fewer claimants than in our model. Later cohorts also live longer, by more at higher incomes \citep{chetty2016association}. \citet{Shoven2014does} find that the value of delaying has risen with the credit, with lower real interest rates, and with longer life expectancy, and that at the low real interest rates of recent years delaying is actuarially advantageous for many individuals, most of all for primary earners in married couples. When the claimant and the government discount at the same rate, such a delay raises the program's outlays (Lemma~\ref{lem:claiming}). We leave a re-evaluation at current-law parameters and current survival to future work.

For the current system, \citet{Jones2023} reach the opposite conclusion about claiming adjustments. In a life-cycle model of claiming, labor supply, and saving in which mortality differs by education, they find that welfare is maximized when baseline benefits do not depend on lifetime earnings and the claiming adjustments are reduced, and that eliminating the earnings test and the income taxation of benefits adds further gains. Three differences between the two exercises bear on the claiming adjustments. First, they search for the welfare-maximizing system starting from current law, under which the credit is 8 percent; we measure the value per dollar of changes from a credit of 4.5 percent. Second, with mortality that differs by education, a larger credit pays more to long-lived high earners, whose delays save the program less (Lemma~\ref{lem:claiming}); in our model, mortality varies with lifetime earnings only through health. Third, in our model part of the effect of a higher credit runs through the earnings test, because benefits withheld after the NRA are credited at the DRC (Section~\ref{sec:fw_claiming}); eliminating the earnings test acts on the same margin directly. Our results therefore concern a different starting point and a different mortality gradient, and they do not settle whether rewards for delay would pay for themselves under current law.

% ---------------------------------------------------------------------------
\subsection{Preferences, Health, and Disability Insurance}\label{sec:disc_prefs}

The reform results rest on the baseline preferences. The logarithmic specification of Appendix~\ref{app:logpref}, with fixed utility costs of participation and reentry, also implies a higher cost of work in poor health than in good health, 0.59 against 0.53 utils, % Source: tables/msm_lnc.tex (FROT tables/estimated_parameters_lnc.tex)
and a large reentry cost. It fits old-age assets better but participation worse, and we do not compute reforms under it.

How much the results depend on preferences differs across families. The sign of the WTP of the DRC increases does not depend on preferences at all: at fixed prices and taxes, an AD reform raises the benefit of every claim after the NRA, including the credit on benefits withheld by the earnings test, and lowers no benefit, so it weakly enlarges every individual's set of feasible choices. Its infinite MVPF depends on preferences only through the responses that determine the sign of its net cost. The structural feature behind these responses is the fixed cost of participation, 791 hours a year in good health and 868 hours in poor health or with a disability (\autoref{tab:estimates}). In the same estimated economy, this nonconvexity is the main driver of nonparticipation at older ages, ahead of the effects of health on wages and on mortality \citep{NdiayeYu2026frot}. What matters for these responses is primarily the level of this cost, which makes retirement a discrete choice and determines how many workers are close to the margin between work and exit at 62 to 69; its health gradient is modest, about 10 percent. % Derived: 867.54/791.07 = 1.097 (tables/est_parameters.tex).
Health sorts workers at this margin through wages, the fixed cost, and access to DI, and, through survival, it determines whether the induced delays occur among claimants for whom the credit is below their actuarially fair rate (Lemma~\ref{lem:claiming}). For the other families, the curvature of utility also governs the WTP, through the value of insurance across health and ability states. Simulated households also hold less wealth at old ages than in the data, likely because the model has no bequest motive (Section~\ref{sec:estimation}), which raises the weight of Social Security in their old-age resources and is likely to raise the value of an additional dollar of benefits.

Disability insurance is an alternative route out of the labor force before the NRA \citep{Yu2024}. Because DI is paid until the NRA without the reduction for early claiming, an NRA increase lengthens DI receipt and returns part of the saving on retirement outlays as DI outlays (Section~\ref{sec:results_claiming}). DI enters the model as an expected benefit, without an application process (Section~\ref{sec:model_di}), so the model captures this extension but not additional applications. \citet{Duggan2007aching} find that the increase in the NRA raised DI enrollment \citep[see also][]{Li2018}; with an application margin, more of the saving from an NRA increase would reappear as DI outlays. The reward-delay reforms leave the NRA, and hence DI eligibility, unchanged.

% ---------------------------------------------------------------------------
\subsection{Distribution}\label{sec:disc_distribution}

Our WTP is ex ante and utilitarian for one cohort: the change in average expected utility at age 30, before health, ability, and survival are realized (Section~\ref{sec:fw_welfare}). It values the insurance a reform provides, but we do not report who gains and who loses by realized health or lifetime earnings. The MVPF framework accommodates such concerns through welfare weights on beneficiaries \citep{hendren2016policy,hendren2020unified,saez2016generalized}, and the rules of each reform indicate how weights would bear on it.

The DRC increases are robust to any positive weights on individuals indexed by their circumstances at age 30: because an AD reform weakly enlarges every individual's set of feasible choices at the prices and taxes of the status quo (Section~\ref{sec:disc_prefs}), no individual is worse off ex ante, so with such weights its WTP is positive and its MVPF infinite. Ex post, after health and wages are realized, an individual can be worse off in some states, because she adjusts saving and work in anticipation of the higher credit, so weights on realized health or lifetime earnings could in principle change the sign. The symmetric reforms lower the benefit of every claim before 65 and raise it after 65, so at unchanged behavior they transfer resources from early to late claimers; in the model, poor health raises the cost of work and lowers survival, both of which favor an early claim (Section~\ref{sec:results_claiming}). In Swedish data, workers who retire earlier consume less in retirement and experience larger drops in consumption at retirement, so that reforms that push workers to retire later carry a substantial consumption-smoothing cost \citep{Kolsrud2024retirement}. Our ex ante WTP values this cost through the marginal utility of consumption in the states in which individuals claim early. The positive WTP of the symmetric reforms is an average over these and other states, and with sufficiently large weights on individuals in poor health it could turn negative. The NRA increases lower the benefit at every claiming age, including for individuals in poor health who cannot offset the cut by working longer; weights that favor them would tend to raise the welfare cost of each dollar saved and to reinforce the advantage of the DRC increases.

Among the other families, weights that favor low lifetime earners would raise the MVPFs of the bottom-bracket expansions relative to the level reforms and lower those of the old-age income-tax cuts, whose gains rise with taxable income in old age, which excludes Social Security benefits (Section~\ref{sec:model_taxes}). Mortality in the model varies with lifetime earnings only through health, whereas in the data life expectancy rises with income \citep{chetty2016association}; with such a gradient, a larger credit would also pay more to long-lived high earners (Section~\ref{sec:disc_cohort}). Section~\ref{sec:financing} discusses the distribution across cohorts.
